% !TeX root = ../thuthesis-example.tex
%
% 第三章：量子网络路由协议综述
%

\chapter{量子网络路由协议综述}
\label{chap:routing-survey}

本章梳理量子网络路由问题的基本约束，介绍代表性量子路由协议的设计思路，并分析其局限性，为后续章节对路由非保序性的理论分析奠定背景。

\section{量子路由问题的基本约束}
\label{sec:routing-constraints}

与经典网络中的最短路径或最优流量路由不同，量子网络路由必须同时应对如下物理约束：

\begin{enumerate}
  \item \textbf{纠缠生成的概率性}：量子信道的Bell对生成成功率$p_e \in (0, 1)$受限于光子传输损耗、探测器效率和光学对准误差等因素，路径的端到端纠缠建立成功率随跳数指数下降。

  \item \textbf{量子态的不可克隆性}：量子信息不能被复制（量子不可克隆定理），因此量子路由无法像经典路由那样通过报文重传或多播分发保证可靠性；一旦纠缠测量失败，该量子资源便不可恢复。

  \item \textbf{保真度衰减}：纠缠对在量子存储器中等待时会因退相干而损失保真度（参见第~\ref{chap:phymodel}章）。在本文的路由度量中，这一效应主要表现为不同链路就绪时间造成的 $T_{\max}$ 对齐等待；在完整仿真中，额外经典信令等待也会由物理层退相干模型反映。

  \item \textbf{存储器相干时间限制}：现有量子中继节点的相干时间$\tau$通常在毫秒至秒级~\cite{lin2024controlflowadaptionefficient}，超过相干时间的等待将使纠缠资源完全退化；路由协议须在有效的相干时间窗口内完成多跳纠缠交换。
\end{enumerate}

上述约束决定了量子路由不能简单将跳数或带宽作为优化目标，而需要联合考虑链路成功率、保真度衰减和相干时间等量子特征参数。

\section{基于保真度乘积的启发式路由}
\label{sec:dijkstra-routing}

一种直觉性的量子路由思路是把路径质量近似写成各链路初始保真度的连乘：若忽略等待退相干和交换过程，端到端保真度可粗略近似为
\[
F_{\rm end}(P) \approx \prod_{e \in P} F^{(0)}_e .
\]
在这个非常粗糙的近似下，可以用 $-\log F^{(0)}_e$ 作为链路权重，把路径选择化约为标准最短路径问题。这个思路的意义主要在于提供一个最小复杂度的基线，而不是作为当前实现中的最终度量。

该近似只有在两个条件同时较强时才勉强成立：

\begin{enumerate}
  \item 各链路的初始保真度大致独立且已知；
  \item 路径内部的等待时间可以忽略，或者所有链路的建立时间都足够接近。
\end{enumerate}

一旦路径内部的建立时间出现明显错配，路径标签就不能再被压缩成一个标量乘积。当前实现采用的正是后一种更完整的路径级语义：整条前缀会在扩展时重新计算 $T_{\max}$，因此“局部最优前缀”并不一定是“接入共享后缀后的全局最优前缀”。这也是后续章节讨论 Dijkstra 基线失效的起点。

\section{当前实现采用的路由基线与切片算法}
\label{sec:current-routing-baseline}

与上述粗略启发式不同，当前 qns-3 的路由层采用的是路径标签级的精确重算语义。每条 label 都保存完整的链路段列表，扩展时重新评估整条路径的 \texttt{predicted\_fidelity} 与 \texttt{t\_max\_ms}，因此路径比较不再是单纯的代价累加，而是“前缀重算 + 时间瓶颈比较”。

在这一实现下，\texttt{DijkstraRoutingProtocol} 的作用是单标签基线。它每次只保留每个节点当前最优的一条 label，因此如果一个前缀因为局部保真度更高而抢占了名额，另一个在当前节点略差、但 \texttt{t\_max\_ms} 更小的前缀就可能被过早剪枝。换句话说，它在当前度量下并不是“经典加性 Dijkstra 的严格最优版本”，而只是 $K=1$ 的 best-first baseline。

与之对应，\texttt{SlicedDijkstraRoutingProtocol} 是多标签搜索。它允许同一节点保留最多 $K$ 条 label，并且可以沿 $T_{\max}$ 轴按固定宽度分桶：
\[
\text{bucket\_id}(p)
=
\left\lfloor \frac{T_{\max}(p)}{\texttt{BucketWidthMs}} \right\rfloor .
\]
当 \texttt{UseBuckets = true} 时，不同 bucket 的 label 可以共存；同一 bucket 内则优先保留 \texttt{predicted\_fidelity} 更高的候选。这样做的目的不是提高单条路径的局部分数，而是尽量保留“快但稍弱”和“慢但更强”的不同前缀切片，从而减少后缀接入后发生排序反转的概率。

当 \texttt{UseBuckets = false} 时，算法退化为基于 \texttt{predicted\_fidelity} 与 \texttt{t\_max\_ms} 的 Pareto 多标签搜索。也就是说，它保留的是关于“端到端保真度更高且时间瓶颈不更差”的一组支配前沿，而不是强行把所有候选压缩成一个标量排名。需要注意的是，切片算法并不是全路径枚举算法；在有限 $K$ 与同桶竞争规则下，它仍可能丢掉未来会变好的前缀，因此本文把它作为缓解非保序剪枝的多标签基线，而不是带严格最优性保证的精确算法。

\section{Caleffi的最优纠缠率路由}
\label{sec:caleffi-routing}

Caleffi~\cite{Caleffi2017Routing}在量子中继网络中建立了端到端\emph{纠缠生成率}（entanglement generation rate）的精确随机模型。该工作将原子-光子纠缠生成概率、光子-光子贝尔态测量不完美性、退相干时间等多种物理因素纳入统一的随机分析框架，推导了任意路径的端到端纠缠率闭合表达式，并以最大化纠缠率为路由优化目标。

Caleffi的核心发现是：纠缠生成率度量具有\emph{非保序性}（non-isotonicity）——在路径拼接时，两条路径的优劣排序可能发生逆转，使得局部最优选择在全局上非最优。这一性质从根本上破坏了标准Dijkstra算法的正确性前提：Dijkstra的贪心选择依赖于路径度量对拼接操作的单调性，而非保序性度量无法满足这一条件。针对该缺陷，Caleffi提出通过路径枚举算法穷举所有候选路径以求得全局最优解。这一工作是量子路由非保序性问题的早期理论奠基之一，但路径枚举的指数级复杂度限制了其在大规模网络中的实用性。

值得注意的是，Caleffi所分析的路由度量是\textbf{纠缠生成率}，与上节保真度乘积方法的优化目标不同；路由非保序性问题在两种度量体系中均以不同形式存在，而下一章将针对本文当前实现中的 \texttt{BottleneckFidelityRoutingMetric} 给出严格的非保序性理论分析。

\section{相关工作}
\label{sec:related-routing}

在上述主流方法之外，若干研究针对特定场景提出了改进方案。Ali等人~\cite{Ali2017Optimal}研究了量子网络编码在多播纠缠分发中的最优性条件，为量子组播路由提供了理论依据。Wang等人~\cite{Wang2025}基于路由代数框架提出了K-最优路径路由（KOP），从代数结构的角度系统分析了量子路由度量的非保序性，并设计了在代数意义上保证收敛性的修正算法。Sutcliffe等人~\cite{Sutcliffe2025}针对多方纠缠分发场景，提出了保真度感知的多路径路由策略，在保证最低保真度阈值的前提下优化分发吞吐量。Shi 和 Qian 提出的 Q-CAST~\cite{shi2020qcast} 则代表了多路径并发和吞吐量优先思路的一个典型分支，可作为本文单标签基线之外的对照背景。

面向卫星量子网络的全球尺度纠缠分发，QuESat~\cite{gu2025quesat}设计了基于低轨道卫星星座的被动光学量子中继网络，并提出了针对星地混合拓扑的路由策略。该工作的出现使得"高延迟-高保真度"卫星链路与"低延迟-低相干时间"地面节点共存的异质场景变得更加现实，对量子路由算法提出了更高的要求。

\section{本章小结}
\label{sec:routing-survey-summary}

现有量子路由协议大致可以分为三类：（1）基于保真度乘积或类似粗粒度度量的启发式最短路径方法；（2）以纠缠率或可达速率为度量的代数分析方法；（3）以多路径并发和吞吐量为中心的协议族。三类方法的共同局限在于：\textbf{都没有把当前实现中的路径级 $T_{\max}$ 重算语义作为核心对象来处理}。因此，单标签 Dijkstra 在当前模型下可能过早剪枝，而多标签切片算法则能够更稳妥地保留不同时间前缀的代表性候选，但它本身仍是有限标签近似。下一章将围绕这一点，给出当前度量的非保序性分析以及 \texttt{Dijkstra} 与 \texttt{Sliced Dijkstra} 的对比证明。
